\subsection{Definition}
\definition{Folge}: (unendliche) Liste von Zahlen, wobei jeder Eintrag ein \textbf{Glied der Folge} genannt wird.\\
Kann auch als Funktion/Abbildung $f:\mathbb{N}_0 \rightarrow \mathbb{R}$ (oder $f:\mathbb{N} \rightarrow \mathbb{R}$) aufgefasst werden.
\newline

\definition{Bilder} $a(n)=a_n$ sind Elemente der Folge.

\subsection{Darstellung}
\definition{direkt}: explizit angeben, wie ein beliebiges Folgenglied berechnet werden kann
$a_n = f(n)$
\newline

\definition{rekursiv}: Konstruktion eines Folgenglieds aus den vorangehenden
$a_n=g(a_{n-1}, a_{n_2},...)$
\newline
$f$ und $g$ nennt man \textbf{Bildungsgesetz}
\newline

\subsection{Konvergenz / Divergenz}
Eine Folge ist \definition{konvergent} falls ein $L \in \mathbb{R}$ existiert, so dass $\forall \epsilon > 0 \exists N > 0$ so dass $\forall n > N : |a_n -L| < \epsilon$.
\newline
Eine konvergente Folge besitzt\\ \theorem{genau einen eindeutigen Grenzwert}.
\newline

$L$ ist der \definition{Grenzwert} $\lim_{n \rightarrow \infty} a_n = L$.
\newline

Falls kein $L$ existiert, nennt man die Folge \definition{divergent}.
\newline